\documentclass[11pt,a4paper]{article}

\usepackage[T1]{fontenc}
\usepackage[utf8]{inputenc}
\usepackage{mathptmx}
\usepackage[scaled=0.9]{helvet}
\usepackage[margin=2.3cm]{geometry}
\usepackage{amsmath,amssymb,amsthm,mathtools}
\usepackage{booktabs}
\usepackage{longtable}
\usepackage{tikz}
\usetikzlibrary{arrows.meta,calc,positioning}
\usepackage{pgfplots}
\pgfplotsset{compat=1.16}
\usepackage{microtype}
\usepackage{enumitem}
\usepackage{tcolorbox}
\usepackage[hidelinks]{hyperref}
\usepackage{fancyhdr}
\usepackage{caption}
\usepackage{colortbl}

\tcbuselibrary{skins,breakable}
\captionsetup{font=small,labelfont=bf}

\renewcommand{\contentsname}{Índice}
\renewcommand{\figurename}{Figura}
\renewcommand{\tablename}{Cuadro}

\pagestyle{fancy}
\fancyhf{}
\fancyhead[L]{\footnotesize Informe de resultados · Vision Rover Challenge}
\fancyhead[R]{\footnotesize\thepage}
\renewcommand{\headrulewidth}{0.4pt}

\definecolor{ok}{RGB}{232,244,234}
\definecolor{mal}{RGB}{250,235,235}
\definecolor{inf}{RGB}{240,240,246}

\newcommand{\grados}{^{\circ}}
\newcommand{\conv}{\operatorname{conv}}
\newcommand{\PASA}{\cellcolor{ok}\textsf{\footnotesize PASA}}
\newcommand{\FALLA}{\cellcolor{mal}\textsf{\footnotesize FALLA}}
\newcommand{\INFO}{\cellcolor{inf}\textsf{\footnotesize INFO}}
\newcommand{\ESPER}{\cellcolor{inf}\textsf{\footnotesize ESPERADO}}

\newtcolorbox{hallazgo}[2]{
  enhanced, breakable, colback=black!3, colframe=black!60,
  boxrule=0.6pt, arc=1pt, left=6pt, right=6pt, top=4pt, bottom=4pt,
  title={\bfseries #1\quad\normalfont #2},
  colbacktitle=black!12, coltitle=black, fonttitle=\bfseries\small}

\newcommand{\campo}[1]{\par\smallskip\noindent\textbf{#1}\quad\ignorespaces}

\begin{document}

\begin{titlepage}
\centering
\vspace*{2cm}
{\Large\scshape Universidad Hispanoamericana}\\[0.2cm]
{\large Ingeniería Mecatrónica}\\[2.2cm]

\rule{\linewidth}{1.2pt}\\[0.4cm]
{\LARGE\bfseries Informe de resultados de la\\[0.18cm]
verificación y validación del modelo matemático}\\[0.4cm]
\rule{\linewidth}{1.2pt}\\[0.8cm]

{\large Ejecución completa del protocolo de pruebas:\\
46 pruebas en cinco niveles, sin hardware físico}\\[2.2cm]

\begin{tabular}{rl}
\textbf{Autor} & Sebastián Vega Amores\\[0.12cm]
\textbf{Documentos de referencia} & \textit{Modelo matemático} (v2) y \textit{Protocolo de pruebas}\\[0.12cm]
\textbf{Fecha} & 2 de septiembre de 2026\\
\end{tabular}

\vfill
\begin{tcolorbox}[colback=black!4,colframe=black!55,boxrule=0.6pt,arc=1pt,width=0.88\linewidth]
\centering\small
\textbf{Veredicto general:} el núcleo geométrico del modelo se sostiene sin excepciones.
Se detectaron \textbf{diez hallazgos}, cuatro de ellos en el repositorio oficial.
Ninguno invalida la conclusión central del modelo, pero el último ---las dimensiones
reales del rover--- reduce su margen operativo en un $60\%$.
\end{tcolorbox}
\end{titlepage}

\tableofcontents
\newpage

% =========================================================================
\section{Resumen ejecutivo}

\subsection{Veredicto}

Se ejecutaron 50 pruebas distribuidas en cinco niveles, desde invariantes algebraicas
hasta misiones completas en lazo cerrado. El resultado global:

\begin{table}[htbp]
\centering
\small
\caption{Resultado por nivel.}
\begin{tabular}{llrrrl}
\toprule
Nivel & Alcance & Pruebas & Pasan & Otras & Observación \\
\midrule
N0 & Consistencia interna & 8 & 8 & --- & 1 defecto de documentación (H-02) \\
N1 & Estimadores vs verdad analítica & 17 & 16 & 1 esperada & la refutación de N1.11 era el objetivo \\
N2 & Visión extremo a extremo & 5 & 4 & 1 informativa & 1 hallazgo mayor (H-03) \\
N3 & Capa de planificación & 12 & 12 & --- & 5 hallazgos (H-04, H-05, H-09, H-11, H-12) \\
N4 & Lazo cerrado & 8 & 2 & 3 informativas & 6 hallazgos (H-06 a H-08, H-13 a H-15) \\
\midrule
\multicolumn{2}{l}{\textbf{Total}} & \textbf{50} & \textbf{42} & \textbf{5} & \\
\bottomrule
\end{tabular}
\end{table}

\begin{tcolorbox}[colback=black!5,colframe=black!65,boxrule=0.7pt,arc=1pt]
\textbf{Lectura de este cuadro.} N0 a N3 ---toda la matemática y toda la geometría de
planificación--- pasa sin excepción. Las tres pruebas que fallan en N4 son misiones de
lazo cerrado, y fallan \emph{después} de incorporar las dimensiones reales del rover
medidas sobre el robot armado. La prueba N4.8 separa las dos causas posibles: con un cubo
por vez el rover entrega $3{,}00$ de $3$; con los tres a la vez, $1{,}00$. \textbf{La
geometría del empuje funciona; lo que falta es una capa de navegación que trate a los
demás cubos como obstáculos.} Esa capa está desarrollada en el modelo y no estaba
implementada en el planificador.
\end{tcolorbox}

\subsection{Las tres afirmaciones centrales del modelo, contrastadas}

\begin{table}[htbp]
\centering
\small
\caption{Contraste de las conclusiones del modelo matemático.}
\begin{tabular}{p{5.6cm}p{4.1cm}p{4.6cm}}
\toprule
Afirmación & Evidencia & Veredicto \\
\midrule
\textbf{El ángulo del cubo se reconstruye en forma cerrada desde una arista o un
vértice más el centro.} &
N1.5, N1.6, N1.8, N1.9: error $\le 1{,}5\times10^{-13}$ grados en $216\,000$ casos. &
\textbf{Confirmado.} Exacto en aritmética de doble precisión. \\
\addlinespace
\textbf{La planificación del empuje no requiere conocer el ángulo}, porque la anchura
proyectada está acotada por $[L,\,L\sqrt2]$. &
N1.12, N1.13, N3.2: cota verificada por dos vías independientes sobre $10^6$ pares
aleatorios más malla de $0{,}1\grados$. N3.10 con el ancho útil real. &
\textbf{Refutado en su consecuencia práctica.} La cota es exacta y el cubo siempre entra,
pero con $W_u = 93{,}50$ el presupuesto de control es negativo en el $22{,}8\%$ de las
orientaciones (H-11): conocer $\theta$ pasa a ser obligatorio. \\
\addlinespace
\textbf{El costo de una ruta lo domina el número de quiebres, no la distancia.} &
N3.6: un quiebre de $90\grados$ equivale a $363\ \mathrm{mm}$ de empuje recto
(el modelo contabiliza $243$, omitiendo el término de reorientación). &
\textbf{Confirmado y reforzado.} El costo real es $49\%$ mayor que el que contabiliza la
fórmula geométrica. \\
\addlinespace
\textbf{El presupuesto lateral de control es de $12{,}57\ \mathrm{mm}$}, derivado del
embudo de tolerancia con el ancho supuesto. &
N4.3: barrido de ruido en lazo cerrado; la ruptura ocurre a $14\ \mathrm{mm}$. N3.2 y
N3.10 con el ancho real. &
\textbf{Numéricamente confirmado, pero el parámetro era falso.} El acuerdo geometría--
simulación fue del $10\%$; con $W_u = 93{,}50$ el presupuesto real va de $11{,}75$ a
$-0{,}68\ \mathrm{mm}$ según la orientación. \\
\bottomrule
\end{tabular}
\end{table}

\subsection{Los quince hallazgos}

\begin{table}[htbp]
\centering
\small
\caption{Hallazgos, ordenados por severidad.}
\begin{tabular}{llp{7.4cm}l}
\toprule
Id & Dónde & Qué & Severidad \\
\midrule
H-11 & Modelo & Con el ancho útil real del canal ($93{,}50\ \mathrm{mm}$) el presupuesto
de control lateral es \textbf{negativo} ($-0{,}68\ \mathrm{mm}$) en el $22{,}8\%$ de las
orientaciones. \textbf{Confirmado por medición.} & \textbf{Crítica} \\
H-12 & Modelo & Las paletas se prolongan $55\ \mathrm{mm}$ por delante y el envolvente
no las contaba: $R_{\mathrm{rov}} = 113{,}5$, no $68{,}4$ ($+66\%$). El anillo de
aproximación ya no libra el cubo. & \textbf{Crítica} \\
H-14 & Planificador & Los otros cubos no eran obstáculos de navegación. Con paletas, el
rover los \emph{engancha} y arrastra media cancha. & \textbf{Alta} \\
H-15 & Planificador & No había retirada antes de girar: el rover barría con la paleta el
cubo que acababa de entregar, sacándolo del depósito. & \textbf{Alta} \\
H-13 & Planificador & Sin garantía de vivacidad: varias ramas devuelven «no te muevas» y
el rover quedaba congelado el resto de la ronda. & Media \\
H-03 & Sistema oficial & La bandera \texttt{confiable} no protege el ángulo:
$57\%$ de falsos negativos con umbral de $5\grados$. & \textbf{Alta} \\
H-01 & Simulador oficial & Radios de empuje y oclusión geométricamente incompatibles
con el cuerpo del rover ($3{,}3\times$ menores). & \textbf{Alta} \\
H-06 & Simulador oficial & Las poses iniciales declaradas hacen que los dos rovers se
solapen físicamente. & \textbf{Alta} \\
H-05 & Modelo & El $36{,}3\%$ de la cancha no admite empuje directo: el multi-tramo es
obligatorio, no opcional. & Media \\
H-07 & Alcance & La cesión por pares no basta: con dos rovers hay colisiones en el
$5{,}5\%$ de las corridas (era el $100\%$ con las dimensiones supuestas). & Media \\
H-04 & Modelo/protocolo & La regla de decisión avance/retroceso correcta es
$A+B>180\grados$, no $A>90\grados$. & Media \\
H-10 & Repositorio & Las dimensiones del rover estaban en el CAD de fabricación y son
otras: canal $93{,}50$, huella $99{,}50\times94{,}00$, no $120\times140$. & \textbf{Alta} \\
H-09 & Modelo & Zona de exclusión de borde: $2{,}8\%$ de la cancha (revisado desde
$11{,}2\%$). El borde es un precipicio, no una pared. & Media \\
H-08 & Contrato & El umbral de latencia de $500\ \mathrm{ms}$ es inerte: la misión ya
falló a $200$. & Media \\
H-02 & Ambos & El rango de $\delta$ es $[-180\grados,180\grados)$, no
$(-180\grados,180\grados]$. & Baja \\
\bottomrule
\end{tabular}
\end{table}

% =========================================================================
\section{Cómo se ejecutó}

\subsection{Bancos de prueba}

Se usaron tres fuentes de verdad, de independencia creciente:

\begin{enumerate}[label=\arabic*.,leftmargin=2.4em,itemsep=3pt]
\item \textbf{Verdad analítica} (N0, N1, N3): la configuración se genera por fórmula y
      el estimador debe invertirla. Detecta errores de implementación y de álgebra.
\item \textbf{Verdad sintética del repositorio oficial} (N2): el generador
      \texttt{vision/sources/generador\_sintetico.py} renderiza imágenes proyectando un
      cubo tridimensional, y el sistema oficial completo ---detección ArUco, homografía,
      pose de cámara, ajuste de silueta--- las procesa. \emph{Ninguna línea del modelo
      participa de la generación}, de modo que la coincidencia es informativa.
\item \textbf{Simulador de dinámica propio} (N4): construido para este protocolo, con la
      geometría de contacto corregida según el hallazgo H-01.
\end{enumerate}

\subsection{Reproducibilidad}

Toda aleatoriedad usa semillas fijas (\texttt{20260902} para los niveles analíticos,
índice de corrida para las misiones). Las corridas son deterministas y repetibles.
El código del banco queda en cinco módulos: \texttt{modelo.py} (implementación de
referencia), \texttt{arnes.py} (registro y criterios), \texttt{n0\_n1.py},
\texttt{n2.py}, \texttt{n3.py}, \texttt{simulador.py}, \texttt{planificador.py} y
\texttt{n4.py}.

\subsection{Entorno}

Python 3 con NumPy, SciPy y OpenCV 4.13 (con \texttt{aruco}). El sistema de visión
oficial se ejecutó sin modificaciones: la prueba N2.1 reproduce su línea base publicada
dentro del $0{,}5\%$, lo que confirma que el banco está correctamente montado antes de
usarlo para medir magnitudes nuevas.

% =========================================================================
\section{Resultados por nivel}

\subsection{N0 --- Consistencia interna}

Ocho pruebas, todas superadas. Los errores máximos están en el orden del épsilon de
máquina, lo que confirma que las convenciones del marco fila-descendente están
correctamente implementadas.

\begin{table}[htbp]
\centering
\small
\caption{Nivel N0. Errores máximos sobre los barridos completos.}
\begin{tabular}{llrl}
\toprule
Id & Prueba & Casos & Error máximo \\
\midrule
N0.1 & Ortogonalidad y determinante de $\mathbf R(\theta)$ & $14\,401$ & $2{,}2\times10^{-16}$ \\
N0.2 & Homomorfismo $\mathbf R(\alpha)\mathbf R(\beta)=\mathbf R(\alpha+\beta)$ & $32\,400$ & $1{,}4\times10^{-15}$ \\
N0.3 & Acción sobre versores y recuperación del argumento & $129\,600$ & $5{,}7\times10^{-14}$ grados \\
N0.4 & Forma cerrada de $\mathbf R(90\grados)$ & $100\,000$ & $6{,}4\times10^{-14}$ \\
N0.5 & Convención fila-descendente (anclada al contrato) & $720$ & $0$ (exacto) \\
N0.6 & Diferencia angular con signo & $667\,440$ & $0$ \\
N0.7 & $\rho_P$ contra minimización explícita & $\sim10^{6}$ & $0$ \\
N0.8 & Identidades del apéndice A & $129\,600$ & $1{,}2\times10^{-15}$ \\
\bottomrule
\end{tabular}
\end{table}

N0.5 merece una nota: se verificó $\mathbf u(\theta)$ contra la fórmula que aparece en
la \emph{sección 4 del contrato oficial} (\texttt{dcol} $=\cos\theta$,
\texttt{drow} $=-\sin\theta$), no contra el modelo. La coincidencia es exacta, bit a
bit. La convención está bien anclada.

\subsection{N1 --- Estimadores contra verdad analítica}

Diecisiete pruebas: dieciséis superadas y una refutación deliberada.

\begin{table}[htbp]
\centering
\small
\caption{Nivel N1. Resultados principales.}
\begin{tabular}{llp{7.0cm}l}
\toprule
Id & Prueba & Resultado & Veredicto \\
\midrule
N1.1 & Propiedades métricas del cuadrado & lado, diagonal, área y centroide correctos a $3\times10^{-14}$ relativo & \PASA \\
N1.2 & Recursión de $90\grados$ & idéntica a la forma trigonométrica ($1{,}1\times10^{-13}$ mm) & \PASA \\
N1.3 & Invariancia $C_4$ & $5{,}7\times10^{-14}$ mm & \PASA \\
N1.4 & Aristas, medios y normales & normales exteriores correctas en los $2\,880$ casos & \PASA \\
N1.5 & B1 (vértice + centro) & $1{,}5\times10^{-13}$ grados en $144\,000$ casos & \PASA \\
N1.6 & B2 (arista) & $1{,}4\times10^{-13}$ grados; invariante al sentido de recorrido & \PASA \\
N1.7 & B2, signo de la normal & $9\,000$ testigos interiores y $19\,440$ realistas: ningún fallo & \PASA \\
N1.8 & B3 (diagonal) & $8{,}5\times10^{-14}$ grados & \PASA \\
N1.9 & B4 sin ruido, 11 subconjuntos & $1{,}1\times10^{-13}$; los diagonales no degeneran & \PASA \\
N1.10 & B4, cota de Cramér--Rao & desvío $\le 0{,}43\%$ de la predicción & \PASA \\
N1.11 & B4 ante vértice atípico & se rompe, como se esperaba & \ESPER \\
N1.12 & Función soporte & forma cerrada $=$ máximo explícito a $6{,}4\times10^{-14}$ & \PASA \\
N1.13 & Cota $[L, L\sqrt2]$ & $10^6$ pares aleatorios, cero contraejemplos & \PASA \\
N1.14 & Lema de silueta & $78\,660$ casos con nadir exterior: todos con 6 vértices & \PASA \\
N1.15 & Paralaje, involución & $1{,}7\times10^{-13}$ mm & \PASA \\
N1.16 & Paralaje, invariancia de dirección & $3{,}4\times10^{-12}$ grados & \PASA \\
N1.17 & Centro proyectivo vs promedio & diagonales exactas; sesgo del promedio caracterizado & \PASA \\
\bottomrule
\end{tabular}
\end{table}

\subsubsection{N1.10 --- La cota de Cramér--Rao se cumple y se conoce su límite}

El barrido en $\sigma$ confirma la fórmula $\sigma_{\hat\theta}=\sigma/(\varrho_{\rm ef}\sqrt n)$
y además delimita su régimen de validez.

\begin{table}[htbp]
\centering
\small
\caption{N1.10. Monte Carlo con $50\,000$ realizaciones por celda,
$L=60\ \mathrm{mm}$, $\theta=31\grados$.}
\begin{tabular}{rrrrrr}
\toprule
$n$ & $\sigma$ [mm] & sesgo [$\grados$] & $\sigma_{\hat\theta}$ empírico & predicho & razón \\
\midrule
2 & $0{,}25$ & $+0{,}0001$ & $0{,}3377$ & $0{,}3376$ & $1{,}0001$ \\
2 & $1{,}00$ & $+0{,}0047$ & $1{,}3467$ & $1{,}3505$ & $0{,}9972$ \\
2 & $4{,}00$ & $-0{,}0094$ & $5{,}4481$ & $5{,}4019$ & $1{,}0086$ \\
3 & $1{,}00$ & $-0{,}0014$ & $0{,}8292$ & $0{,}8270$ & $1{,}0026$ \\
4 & $0{,}25$ & $-0{,}0006$ & $0{,}1693$ & $0{,}1688$ & $1{,}0028$ \\
4 & $1{,}00$ & $-0{,}0004$ & $0{,}6752$ & $0{,}6752$ & $0{,}9999$ \\
4 & $2{,}00$ & $-0{,}0105$ & $1{,}3530$ & $1{,}3505$ & $1{,}0019$ \\
4 & $4{,}00$ & $-0{,}0356$ & $2{,}7189$ & $2{,}7009$ & $1{,}0067$ \\
\bottomrule
\end{tabular}
\end{table}

El estimador es insesgado y eficiente hasta $\sigma = 2\ \mathrm{mm}$ con desvío inferior
al $0{,}5\%$. A $\sigma = 4\ \mathrm{mm}$ aparece un sesgo pequeño pero estadísticamente
detectable ($z = 2{,}9$ con $n=4$), lo que importa porque el peor error medido del
sistema de visión bajo oclusión es $4{,}88\ \mathrm{mm}$: se está justo en el borde del
régimen donde la aproximación de primer orden empieza a ceder.

\subsubsection{N1.11 --- La refutación buscada}

Esta prueba se diseñó para reprobar, y reprobó. Contaminando un vértice con un
desplazamiento $d$:

\begin{table}[htbp]
\centering
\small
\caption{N1.11. Sensibilidad de B4 a un vértice atípico ($n=4$).}
\begin{tabular}{rrrr}
\toprule
$d$ [mm] & error de $\hat\theta$ [$\grados$] & error de $\hat{\mathbf c}$ [mm] & pendiente [$\grados$/mm] \\
\midrule
$5$ & $1{,}69$ & $1{,}25$ & $0{,}338$ \\
$20$ & $6{,}73$ & $5{,}00$ & $0{,}336$ \\
$80$ & $28{,}12$ & $20{,}00$ & $0{,}362$ \\
$120$ & $44{,}99$ & $30{,}00$ & $0{,}422$ \\
$200$ & $41{,}66$ & $50{,}00$ & --- \\
\bottomrule
\end{tabular}
\end{table}

El error del centro es exactamente $d/n$, como corresponde a una media aritmética: el
estimador de mínimos cuadrados \textbf{no tiene punto de ruptura}. El error angular
parece saturarse a $200\ \mathrm{mm}$, pero eso es un artefacto de la métrica: $\rho_{90}$
no puede superar $45\grados$. El error real sigue creciendo sin cota.

\begin{tcolorbox}[colback=black!3,colframe=black!50,boxrule=0.5pt,arc=1pt]
\textbf{Consecuencia de diseño.} B4 (Procrustes) es el estimador \emph{óptimo} cuando
todos los vértices son del cubo, y el \emph{peor} cuando uno no lo es. Como un extractor
de vértices real producirá falsos positivos ---una esquina del rover, un reflejo---,
B4 no puede usarse solo. La cascada correcta es: estimar $\theta$ por B2 sobre
\emph{cada} arista candidata, tomar la mediana circular, descartar los vértices
incompatibles con esa mediana, y recién entonces refinar con B4 sobre los
supervivientes.
\end{tcolorbox}

\subsubsection{N1.14 --- Las degeneraciones del lema de silueta}

El lema afirma «genéricamente seis vértices». La prueba buscó activamente las
configuraciones no genéricas:

\begin{table}[htbp]
\centering
\small
\caption{N1.14. Conteo de vértices de la envolvente convexa.}
\begin{tabular}{lrl}
\toprule
Posición del nadir & Casos & Vértices del casco \\
\midrule
Exterior a la base (malla de $2\times2$ m, tres alturas de cámara) & $78\,660$ & 6 en el $100\%$ \\
Interior a la base & $60$ & 4 en el $100\%$ \\
Sobre un vértice de la base & $24$ & 4 \\
Sobre el punto medio de una arista & $24$ & 4 en 16 casos, 5 en 8 \\
\bottomrule
\end{tabular}
\end{table}

El lema queda confirmado y sus degeneraciones caracterizadas: el nadir sobre un vértice
produce $B \subset T$ (casco de cuatro vértices, porque escalar un cuadrado alrededor de
uno de sus propios vértices lo contiene), y sobre el punto medio de una arista produce
cinco. Ambas son configuraciones de medida nula que no se dan en la práctica, pero ahora
están documentadas en lugar de supuestas.

\subsection{N2 --- Visión extremo a extremo}

Este es el nivel que aportó el hallazgo más importante.

\subsubsection{N2.1 --- La línea base se reproduce}

\begin{table}[htbp]
\centering
\small
\caption{N2.1. Error de posición del detector oficial, en milímetros.}
\begin{tabular}{lrrrl}
\toprule
Escenario & cenital & inclinada & residuo & documentado \\
\midrule
Los tres cubos de la configuración & $1{,}23$ & $0{,}76$ & $0{,}017$ & --- \\
Cubos repartidos por la cancha & $1{,}40$ & $\mathbf{1{,}05}$ & $0{,}015$ & $1{,}05$ \\
Rotaciones de $0\grados$, $22{,}5\grados$, $45\grados$ & $0{,}84$ & $0{,}65$ & $0{,}016$ & --- \\
Un rover empujando el cubo & $1{,}40$ & $\mathbf{4{,}88}$ & $0{,}133$ & $4{,}88$ \\
Un rover tapando más del cubo & $34{,}31$ & $12{,}63$ & $0{,}244$ & no confiable \\
\bottomrule
\end{tabular}
\end{table}

Desvío respecto de los valores publicados: $0{,}47\%$ y $0{,}10\%$. El banco está bien
montado.

\subsubsection{N2.2 --- Primera caracterización del error angular}

Ninguna herramienta del repositorio mide esta magnitud. Barriendo
$\theta\in[0\grados,90\grados)$ con paso $2{,}5\grados$ sobre cinco zonas de la cancha
y los dos modos de cámara ($360$ observaciones):

\begin{table}[htbp]
\centering
\small
\caption{N2.2. Error angular del detector oficial, cubo despejado.}
\begin{tabular}{lrrrr}
\toprule
Modo de cámara & $n$ & medio [$\grados$] & p95 [$\grados$] & máximo [$\grados$] \\
\midrule
Cenital & $180$ & $0{,}419$ & $1{,}00$ & $1{,}50$ \\
Inclinada $8\grados$ & $180$ & $0{,}433$ & $1{,}00$ & $1{,}50$ \\
\bottomrule
\end{tabular}
\end{table}

El error no depende de $\theta$ ni de la zona, y está cuantizado en múltiplos de
$0{,}25\grados$, que es exactamente la resolución final del descenso por coordenadas
($4\grados/2^4$). Es decir: \textbf{el error angular con el cubo despejado está
dominado por la discretización del optimizador, no por el ruido de la imagen.} Bajar
ese error es cuestión de agregar una ronda de refinamiento, no de mejorar la cámara.

\subsubsection{N2.3 --- Y su degradación bajo oclusión}

\begin{table}[htbp]
\centering
\small
\caption{N2.3. Degradación con oclusión (cámara inclinada, $18$ ángulos por condición).}
\begin{tabular}{lrrrr}
\toprule
Condición & pos.\ máx.\ [mm] & áng.\ medio [$\grados$] & áng.\ máx.\ [$\grados$] & frac.\ confiable \\
\midrule
Despejado & $0{,}90$ & $0{,}39$ & $1{,}00$ & $1{,}00$ \\
Rover empujando ($\approx22\%$) & $8{,}59$ & $\mathbf{10{,}92}$ & $\mathbf{19{,}00}$ & $\mathbf{1{,}00}$ \\
Rover tapando ($\approx70\%$) & $43{,}89$ & $21{,}94$ & $43{,}50$ & $0{,}28$ \\
\bottomrule
\end{tabular}
\end{table}

Con un rover empujando, la posición se degrada $9{,}5\times$ (de $0{,}90$ a
$8{,}59\ \mathrm{mm}$) pero el ángulo se degrada $\mathbf{28\times}$ (de $0{,}39$ a
$10{,}92\grados$ de media). Y \textbf{el $100\%$ de esas detecciones se declara
confiable}. De aquí sale el hallazgo H-03.

\subsubsection{N2.5 --- El eslabón débil de la familia analítica}

Extrayendo vértices de la base del contorno real por el criterio de homotecia
($\mathbf q = \mathbf N + \kappa(\mathbf p - \mathbf N)$), sobre $90$ escenas:

\begin{table}[htbp]
\centering
\small
\caption{N2.5. Extracción de vértices de la base y precisión resultante.}
\begin{tabular}{lr}
\toprule
Escenas con 0 vértices de base recuperados & $32$ ($36\%$) \\
Escenas con 1 vértice & $20$ ($22\%$) \\
Escenas con 2 vértices & $38$ ($42\%$) \\
Escenas con una \emph{arista} completa de base & $22$ ($24\%$) \\
\midrule
B2 sobre arista real: error medio / máximo & $0{,}340\grados$ / $0{,}897\grados$ \\
Ajuste oficial de silueta: error medio / máximo & $0{,}456\grados$ / $1{,}500\grados$ \\
\bottomrule
\end{tabular}
\end{table}

El resultado es de doble filo y ambos filos importan. Por un lado, \textbf{cuando la
arista se recupera, B2 es más preciso que el ajuste oficial} ($0{,}34\grados$ contra
$0{,}46$ de media, y $0{,}90$ contra $1{,}50$ de máximo): la fórmula cerrada no tiene
error de discretización. Por otro, la arista solo se recupera en el $24\%$ de las
escenas. \emph{El eslabón débil de la familia analítica no es la matemática sino la
extracción de vértices}, que es justamente el paso que el modelo no desarrolla. La
conclusión práctica es que B2 sirve como \textbf{refinador} de una pose ya estimada por
el ajuste de silueta, no como estimador primario.

\subsection{N3 --- Capa de planificación}

Ocho pruebas, todas superadas, con dos hallazgos.

\begin{table}[htbp]
\centering
\small
\caption{Nivel N3. Resultados.}
\begin{tabular}{llp{7.4cm}l}
\toprule
Id & Prueba & Resultado & Veredicto \\
\midrule
N3.1 & Mapa de factibilidad & $36{,}3\%$ de la cancha sin empuje directo (H-05) & \PASA \\
N3.2 & Embudo por dos vías & coinciden a $3{,}6\times10^{-14}$ mm; cota universal & \PASA \\
N3.3 & Contención del corredor & $10^4$ trayectos con rotación $\pm30\grados$: cero violaciones & \PASA \\
N3.4 & Funcional \textsc{gtg} & mínimo correcto; regla de decisión corregida (H-04) & \PASA \\
N3.5 & Umbral de cuerda del anillo & teórico $84{,}95\grados$, medido $85{,}0\grados$ & \PASA \\
N3.6 & Costo de un quiebre & $399\ \mathrm{mm}$ equivalentes para $90\grados$ & \PASA \\
N3.7 & Conjunto alcanzable & cero falsos negativos en $10^4$ configuraciones & \PASA \\
N3.8 & Interferencia entre corredores & $46\%$ de los pares interfiere & \PASA \\
\bottomrule
\end{tabular}
\end{table}

\subsubsection{N3.6 --- El costo de doblar, cuantificado}

\begin{table}[htbp]
\centering
\small
\caption{N3.6. Costo de un quiebre expresado en distancia rectilínea equivalente,
con $v=120\ \mathrm{mm/s}$ y $\omega=90\grados/\mathrm{s}$.}
\begin{tabular}{rrrr}
\toprule
Quiebre & arco [mm] & giro [s] & equivalente [mm] \\
\midrule
$15\grados$ & $39{,}9$ & $0{,}17$ & $99{,}9$ \\
$45\grados$ & $119{,}7$ & $0{,}50$ & $219{,}7$ \\
$90\grados$ & $239{,}4$ & $1{,}00$ & $\mathbf{399{,}4}$ \\
$180\grados$ & $478{,}9$ & $2{,}00$ & $758{,}9$ \\
\bottomrule
\end{tabular}
\end{table}

El modelo afirmaba $\approx279\ \mathrm{mm}$ para un quiebre de $90\grados$. Ese número
es correcto para el término retirada $+$ arco ($279{,}4\ \mathrm{mm}$ medidos, desvío
$0{,}01\%$), pero \textbf{omite el tiempo de reorientación}. Incluyéndolo, el costo real
es $399\ \mathrm{mm}$: un $43\%$ más. La regla de diseño «minimizar quiebres antes que
longitud» sale reforzada, no debilitada.

\subsubsection{N3.8 --- La coordinación no es opcional}

Con corredores de empuje aleatorios, el $\mathbf{46\%}$ de los pares interfiere ---sus
corredores distan menos de $L\sqrt2 = 84{,}85\ \mathrm{mm}$---. La condición analítica
coincide con el muestreo denso en el $99{,}99\%$ de $10^4$ casos. Es decir: si se asignan
tareas a los dos rovers sin mirar la geometría, casi la mitad de las asignaciones son
incompatibles. Esto anticipa el hallazgo H-07.

\subsection{N4 --- Lazo cerrado}

Se construyó el simulador de dinámica especificado en la sección 5.1 del protocolo:
cinemática de uniciclo, contacto entre cuerpos rígidos por el teorema del eje separador,
empuje direccional restringido a la cara delantera, rotación inducida por centro
instantáneo cuasiestático, oclusión por área tapada y telemetría v1 con ruido, pérdidas
y latencia configurables.

Las pruebas se corrieron en dos configuraciones: \textbf{un rover}, que aísla la capa
geométrica del modelo (su alcance declarado), y \textbf{dos rovers}, que expone la falta
de la capa de coordinación multi-agente (explícitamente fuera de alcance).

\subsubsection{Una corrección al propio simulador}

Durante la puesta a punto se detectó que el primer modelo de oclusión del simulador
---superposición de la \emph{huella} del rover con la \emph{base} del cubo--- no puede
producir oclusión durante un empuje: vistos desde arriba, los dos cuerpos se tocan pero
no se solapan. Es exactamente el mismo defecto que se le señala al simulador oficial en
el hallazgo H-01.

La corrección consistió en comparar \textbf{siluetas} en lugar de huellas. La silueta de
un cuerpo de altura $h$ es $\conv(B \cup T)$ con $T = \mathbf N + \kappa(h)(B-\mathbf N)$,
donde $\mathbf N$ es el nadir. Como el rover mide $90\ \mathrm{mm}$ y el cubo
$60\ \mathrm{mm}$, la silueta del rover se proyecta sobre la del cubo, que es el
mecanismo físico real de la oclusión. Con el modelo corregido la oclusión aparece donde
debe:

\begin{table}[htbp]
\centering
\small
\caption{Oclusión medida en el simulador corregido, con el rover en pose de contacto.}
\begin{tabular}{lrrr}
\toprule
Posición del cubo & dist.\ al nadir [celdas] & $d=5{,}00$ c & $d=5{,}62$ c \\
\midrule
$(21{,}5;\,21{,}5)$ --- bajo la cámara & $0$ & $3{,}2\%$ & $0{,}0\%$ \\
$(26;\,10)$ & $12{,}4$ & $5{,}5\%$ & $0{,}0\%$ \\
$(10;\,10)$ & $16{,}3$ & $9{,}5\%$ & $0{,}2\%$ \\
$(34;\,9)$ & $17{,}7$ & $13{,}8\%$ & $0{,}7\%$ \\
\bottomrule
\end{tabular}
\end{table}

La estructura es la correcta: bajo la cámara casi no hay oclusión, y crece al alejarse
del nadir. La magnitud queda por debajo del $22\%$ que reporta el sistema oficial,
porque el generador oficial dibuja además las caras laterales del rover; el mecanismo,
sin embargo, es el mismo. \textbf{Se deja registrado que el autor cometió el mismo error
que estaba auditando}, porque es precisamente la clase de defecto que solo aparece
cuando se construye el escenario y se lo mira, y no cuando se razona sobre él.

\subsubsection{Resultados}

\begin{table}[htbp]
\centering
\small
\caption{N4.1 y N4.2. Misión nominal. Las dos configuraciones usan el mismo
planificador; la única diferencia es el número de rovers.}
\begin{tabular}{llrrrrr}
\toprule
Prueba & Config & entregas & $t_{50}$ [s] & $t_{90}$ [s] & colisiones & corridas limpias \\
\midrule
N4.1 (sin ruido) & 1 rover  & $3/3$ & $33{,}4$ & --- & $0$ & --- \\
N4.1 (sin ruido) & 2 rovers & $3/3$ & $21{,}7$ & --- & $\mathbf{0}$ & --- \\
\addlinespace
N4.2 (200 semillas) & 1 rover  & $\mathbf{98{,}5\%}$ & $43{,}6$ & $48{,}4$ & $0{,}0$ & $\mathbf{100\%}$ \\
N4.2 (200 semillas) & 2 rovers & $96{,}5\%$ & $28{,}4$ & $32{,}6$ & $0{,}2$ & $\mathbf{94{,}5\%}$ \\
\bottomrule
\end{tabular}
\end{table}

Dos rovers son un $35\%$ más rápidos en la mediana ($28{,}4$ contra $43{,}6\ \mathrm{s}$)
y colisionan en el $5{,}5\%$ de las corridas. Ese es el intercambio que cuantifica el
hallazgo H-07: el paralelismo paga, y la falta de coordinación cuesta --- aunque
\textbf{mucho menos de lo que costaba} con las dimensiones supuestas, donde colisionaban
en el $100\%$ de las corridas. Es el efecto colateral benigno del rover real: al ser más
chico, dos cuerpos se estorban mucho menos en la misma cancha.

\subsubsection{N4.3 --- El presupuesto de error del modelo, falsado}

Esta es la prueba que convierte el modelo en algo refutable. El modelo predice un
presupuesto lateral de control de $e^{\max}_{\mathrm{control}} = 12{,}57\ \mathrm{mm}$.

\begin{table}[htbp]
\centering
\small
\caption{N4.3. Tasa de entrega completa contra ruido posicional, un rover,
40 semillas por punto.}
\begin{tabular}{rrrrrrrrrr}
\toprule
$\sigma$ [mm] & $0$ & $1$ & $2$ & $4$ & $6$ & $8$ & $10$ & $\mathbf{14}$ & $20$ \\
\midrule
tasa $3/3$ & $1{,}00$ & $1{,}00$ & $0{,}93$ & $0{,}93$ & $0{,}83$ & $0{,}65$ & $0{,}63$ &
$\mathbf{0{,}33}$ & $0{,}08$ \\
\bottomrule
\end{tabular}
\end{table}

El punto de ruptura (tasa $<0{,}5$) está en $\sigma = 14\ \mathrm{mm}$, es decir
$\mathbf{1{,}1\times}$ \textbf{el presupuesto predicho}. Para una predicción derivada de
geometría pura ---la anchura proyectada de un cuadrado contra el ancho de las paletas,
menos la incertidumbre de la visión--- contrastada contra una simulación dinámica
completa con contacto entre cuerpos, ruido, oclusión y realimentación, un acuerdo del
$10\%$ es un resultado fuerte. \textbf{El embudo de tolerancia del modelo predice
correctamente dónde se rompe el sistema.}

\subsubsection{N4.4 --- La latencia}

\begin{table}[htbp]
\centering
\small
\caption{N4.4. Tasa de entrega completa contra latencia, un rover. Las dos columnas
corresponden a dos presupuestos de tiempo de misión distintos.}
\begin{tabular}{rrrrrrrr}
\toprule
Latencia [ms] & $0$ & $50$ & $100$ & $200$ & $400$ & $800$ & $1600$ \\
\midrule
tasa con $t_{\max}=180\ \mathrm{s}$ & $0{,}97$ & $0{,}97$ & $1{,}00$ & \multicolumn{4}{c}{---} \\
tasa con $t_{\max}=60\ \mathrm{s}$ & $1{,}00$ & $1{,}00$ & $\mathbf{0{,}17}$ & $0{,}00$ & $0{,}00$ & $0{,}00$ & $0{,}00$ \\
frenadas por la regla RI4 & $0$ & $0$ & $0$ & $0$ & $0$ & $14\,220$ & $14\,028$ \\
\bottomrule
\end{tabular}
\end{table}

Las dos filas dicen cosas distintas y complementarias. Con tiempo de sobra, el
planificador \emph{tolera} $100\ \mathrm{ms}$ de latencia: entrega igual. Con un
presupuesto de $60\ \mathrm{s}$, a $100\ \mathrm{ms}$ solo el $17\%$ de las corridas
llega a tiempo. Es decir: \textbf{la latencia no rompe la corrección, rompe el tiempo}.
Y a partir de $200\ \mathrm{ms}$ falla en ambos sentidos.

Las columnas con y sin la regla de corte RI4 son idénticas en todo el rango, porque la
regla solo dispara por encima de $500\ \mathrm{ms}$, donde ya no queda nada que salvar.
De ahí el hallazgo H-08.

\subsubsection{N4.5 --- Escenarios adversarios}

\begin{table}[htbp]
\centering
\small
\caption{N4.5. Seis configuraciones construidas para romper hipótesis nombradas.}
\begin{tabular}{lrrrrl}
\toprule
Escenario & entregas & $t$ [s] & colis. & salidas & desenlace \\
\midrule
(a) cubo en la esquina opuesta & $0/1$ & $120$ & $0$ & $0$ & abandonado: zona trampa \\
(b) cubo pegado a la pared & $1/1$ & $17{,}7$ & $0$ & $0$ & entregado \\
(c) dos cubos adyacentes & $2/2$ & $28{,}5$ & $0$ & $0$ & entregados \\
(d) cubo ajeno en el corredor & $1/2$ & $120$ & $0$ & $0$ & el ajeno cayó en la trampa \\
(e) corredores cruzados & $1/2$ & $120$ & $5$ & $0$ & colisiones (H-07) \\
(f) cubo sobre depósito ajeno & $0/1$ & $120$ & $0$ & $0$ & abandonado: zona trampa \\
\bottomrule
\end{tabular}
\end{table}

\textbf{Ningún rover salió de la cancha en ningún escenario.} Los casos (a) y (f) no son
fallas del planificador: colocan el cubo en la zona trampa del hallazgo H-09, donde no
existe secuencia de empujes que lo saque. Lo que la prueba verifica es que el
planificador \emph{lo detecta y abandona}, en vez de gastar la ronda empujando contra
una pared. El caso (d) es más interesante: el cubo azul, que empezó en el corredor,
terminó empujado contra la pared por el tránsito del verde y quedó atrapado. Es la
demostración dinámica de por qué la regla «nunca empujar hacia una pared» tiene que
aplicarse también a los cubos que uno \emph{no} está transportando.

\subsubsection{N4.6 --- Oclusión durante el transporte}

Con el modelo de siluetas corregido, sobre $121$ muestras de la fase de empuje:

\begin{table}[htbp]
\centering
\small
\caption{N4.6. Oclusión y error lateral durante un empuje completo.}
\begin{tabular}{lrr}
\toprule
Magnitud & media & máximo \\
\midrule
Fracción de la silueta del cubo tapada & $6{,}2\%$ & $26{,}6\%$ \\
\texttt{age\_ms} del cubo & $0$ & $\mathbf{0}$ \\
Error lateral $|e|$ [mm] & $5{,}07$ & $13{,}37$ \\
\bottomrule
\end{tabular}
\end{table}

El resultado corrige una expectativa del modelo. La oclusión durante un empuje normal
llega al $26{,}6\%$ ---coherente con el $22\%$ que reporta el sistema oficial--- pero
\textbf{no alcanza para perder la detección}: \texttt{age\_ms} se mantiene en cero
durante todo el transporte. La consecuencia es importante y va en contra de la
intuición:

\begin{tcolorbox}[colback=black!3,colframe=black!50,boxrule=0.5pt,arc=1pt]
\textbf{Durante el empuje el dato no llega viejo: llega fresco y equivocado.} El riesgo
no es la obsolescencia sino la \emph{degradación silenciosa de la exactitud}: posición
con $4{,}88\ \mathrm{mm}$ de error en vez de $1{,}05$, y ángulo con hasta $19\grados$ en
vez de $1{,}5$ ---todo con \texttt{age\_ms} $=0$ y la bandera \texttt{confiable} en
verdadero (H-03). Una lógica que solo vigile \texttt{age\_ms} no ve nada de esto.
\end{tcolorbox}

El error lateral máximo, $13{,}37\ \mathrm{mm}$, supera levemente el presupuesto de
$12{,}57$ y aun así la entrega se completó: el presupuesto del modelo es conservador,
porque supone simultáneamente el peor ángulo del cubo y el peor error de visión.

\newpage
% =========================================================================
\section{Los hallazgos en detalle}

\begin{hallazgo}{H-11}{El presupuesto de control lateral es negativo en el $22{,}8\%$ de las orientaciones}
\campo{Dónde} Consecuencia del ancho útil real del canal, establecido sobre
\texttt{CENFOBOT\_PIEZAS\_SEPARADAS.dxf}. Verificado en N3.2 y en la prueba nueva N3.10.
\campo{Qué} El chasis mide $93{,}50\ \mathrm{mm}$ de cuerpo, con lengüetas pasantes de
$3{,}00\ \mathrm{mm}$ por lado. Como las lengüetas atraviesan los paneles laterales, las
\emph{caras interiores} de los paneles coinciden con los bordes del cuerpo del chasis, y
el ancho útil del canal es
\[
  W_u = 93{,}50\ \mathrm{mm}, \qquad W_{\mathrm{ext}} = 99{,}50\ \mathrm{mm}.
\]
Con $L = 60$, la tolerancia lateral cae a
$e_{\max} \in [\,4{,}32;\ 16{,}75\,]\ \mathrm{mm}$. Descontando la incertidumbre de la
visión ($\sigma_{\mathrm{vis}} = 5\ \mathrm{mm}$), el presupuesto para el error de
control en el peor caso es
\[
  4{,}32 - 5{,}00 \;=\; \mathbf{-0{,}68\ \mathrm{mm}} .
\]
Es negativo: la sola incertidumbre de la medición excede la tolerancia disponible, aun
con un controlador perfecto. La banda afectada es
$\Delta \in (34{,}76\grados;\ 55{,}24\grados) \bmod 90\grados$, el $22{,}8\%$ de las
orientaciones posibles.
\campo{Impacto} \textbf{Crítico, y cambia el alcance del proyecto.} En las dos revisiones
anteriores conocer el ángulo del cubo era, primero, una mejora de eficiencia, y luego una
palanca de robustez. Con el ancho real \textbf{es condición de funcionamiento}: existen
orientaciones para las cuales el empuje ciego no tiene margen alguno, y ocurren en casi
una de cada cuatro presentaciones aleatorias del cubo.
\campo{Acción} Regla de apuntado, verificada en N3.10:
\begin{enumerate}[label=(\alph*),leftmargin=2em,itemsep=1pt,topsep=2pt]
\item \textbf{Apuntar el empuje a menos de $21\grados$ de la normal a una cara del cubo}
      (ventana con $3\ \mathrm{mm}$ de margen de control). Como las caras se repiten cada
      $90\grados$, hay cuatro direcciones admisibles y la restricción cuesta poco en
      planificación.
\item \textbf{Implementar la estimación de $\theta$}, que el detector oficial calcula
      internamente pero el contrato v1 no publica. Los estimadores cerrados de la sección
      de resultados N1 la resuelven con una arista o con un vértice más el centro.
\item \textbf{Reevaluar la alineación durante el transporte.} El cubo gira libremente
      dentro del canal ($W_u > L\sqrt2$), así que puede entrar alineado y salir de la
      ventana sin que nada lo impida.
\item \textbf{Verificar con calibre} que el robot armado mide $99{,}50\ \mathrm{mm}$
      punta a punta. Es la única comprobación pendiente, y confirma o refuta todo lo
      anterior en un minuto.
\end{enumerate}
\end{hallazgo}

\begin{hallazgo}{H-12}{El envolvente del rover estaba subestimado un 66\,\%}
\campo{Dónde} Consecuencia de medir el alcance de las paletas sobre el rover armado
($\lambda = 55\ \mathrm{mm}$, medición M6). Verificado en N3.12.
\campo{Qué} El modelo tomaba $R_{\mathrm{rov}} = \tfrac12\sqrt{W_{\mathrm{ext}}^2+\Lambda^2}
= 68{,}44\ \mathrm{mm}$, del contorno del chasis. Pero las paletas se prolongan
$55\ \mathrm{mm}$ por delante de la placa frontal, de modo que la huella \textbf{no es
simétrica} respecto del centro del chasis: llega $47\ \mathrm{mm}$ hacia atrás y
$102\ \mathrm{mm}$ hacia adelante. El disco envolvente mide
\[
  R_{\mathrm{rov}} = \sqrt{102^2 + 49{,}75^2} = \mathbf{113{,}49\ \mathrm{mm}},
\]
un $66\%$ más. La medición del largo total lo confirma por otra vía: $94 + 55 = 149$
contra los $150\ \mathrm{mm}$ medidos.
\campo{Impacto} \textbf{Crítico, y del lado peligroso.} El planificador creía pasar por
huecos por los que no pasa. Tres consecuencias:
\begin{enumerate}[label=(\alph*),leftmargin=2em,itemsep=1pt,topsep=2pt]
\item \textbf{El anillo de aproximación deja de librar el cubo.} Rodear el cubo a
      $\varrho_{\mathrm{puesta}} = 129{,}43\ \mathrm{mm}$ lo \emph{barre} con las paletas.
      Girar detrás de él exige $\varrho_{\mathrm{man}} = 155{,}9\ \mathrm{mm}$.
\item \textbf{La zona trampa crece.} De $5{,}03\%$ a $6{,}46\%$ de la cancha, y la banda
      junto al borde de $2{,}97$ a $4{,}55$ celdas. Los depósitos están a $2{,}5$ celdas
      del borde: quedan \emph{dentro} de esa banda.
\item \textbf{La distancia de contacto no cambia.} Como $\lambda > L\sqrt2/2$, el cubo
      entra en el canal hasta tocar la placa frontal, que es la que empuja.
\end{enumerate}
\campo{Acción} Acceso en dos tramos: al anillo de maniobra primero, y desde ahí en línea
recta por el propio eje de empuje. Implementado y verificado en N4.
\end{hallazgo}

\begin{hallazgo}{H-14}{Los demás cubos no eran obstáculos de navegación}
\campo{Dónde} Planificador. Descubierto al incorporar la forma de U al simulador.
\campo{Qué} El planificador navegaba ignorando los cubos que no eran su objetivo. Con la
huella rectangular anterior eso era casi inocuo: pasar cerca de un cubo lo rozaba. Con las
paletas, la punta \textbf{engancha} el cubo y lo arrastra: en una corrida nominal un cubo
fue desplazado $36$ celdas respecto de su depósito.
\campo{Impacto} \textbf{Alto.} El despeje lateral que el rover necesita al pasar junto a
un cubo ajeno es $W_{\mathrm{ext}}/2 + L\sqrt2/2 = 92{,}2\ \mathrm{mm} = 4{,}6$ celdas, y
el corredor de empuje también debe respetarlo. Ninguna de las dos condiciones se
verificaba.
\campo{Acción} Tratar los demás cubos como obstáculos, tanto en la navegación como en el
corredor de empuje. El modelo ya tenía la maquinaria ---inflado de Minkowski, grafo de
visibilidad, condición de corredores disjuntos de N3.8--- pero el planificador nunca la
usó porque con la huella supuesta no hacía falta.
\end{hallazgo}

\begin{hallazgo}{H-15}{El rover deshacía su propio trabajo al girar}
\campo{Dónde} Planificador. Cuantificado en N4.8.
\campo{Qué} Al terminar un empuje, el rover soltaba el objetivo y giraba en el lugar para
ir al siguiente. Con las paletas por delante, ese giro \textbf{barre el cubo que acaba de
dejar} y lo saca del depósito. En la misión nominal con tres cubos se contabilizan
$\mathbf{4{,}00}$ des-entregas por corrida: cubos que llegaron a estar dentro de la
tolerancia de entrega y después salieron.
\campo{Impacto} \textbf{Alto.} No es que el rover falle en entregar: es que deshace
entregas ya conseguidas. Es el peor tipo de falla porque el trabajo perdido ya estaba
pagado.
\campo{Acción} Retirada recta antes de girar. Para que las puntas queden libres del cubo
hay que retroceder exactamente $\lambda = 55\ \mathrm{mm}$: el rover tiene que desandar el
canal antes de darse vuelta. El costo ya estaba contemplado en el funcional GTG del modelo
---el término de retirada--- pero el planificador nunca lo ejecutaba.
\end{hallazgo}

\begin{hallazgo}{H-13}{El planificador no tenía garantía de vivacidad}
\campo{Dónde} Planificador, varias ramas de la función de decisión.
\campo{Qué} Varias ramas devuelven «no te muevas» ($v=0$, $\omega=0$). Si la condición que
las dispara no cambia por sí sola, el rover queda inmóvil el resto de la ronda. Con la
huella chica el caso no se alcanzaba. Al incorporar las paletas, la rama «sin pose de
puesta válida, sin waypoint, y a menos de $3$ celdas del depósito» pasó a ser alcanzable:
el rover se congelaba con dos cubos sin entregar y $200$ eventos de borde por corrida.
\campo{Impacto} Medio, pero insidioso: un bloqueo silencioso no se distingue de un rover
que está pensando.
\campo{Acción} Vigilante por síntoma ---objetivo fijo y rover quieto más de $2{,}5$
segundos--- que libera el objetivo. No repara la causa; garantiza que la misión siga
avanzando. Bajó los eventos de borde de $200$ a $5$ por corrida.
\end{hallazgo}

\begin{hallazgo}{H-01}{El modelo de contacto del simulador oficial es geométricamente imposible}
\campo{Dónde} \texttt{contrato/config\_simulador.json}.
\campo{Qué} Declara \texttt{radio\_empuje\_celdas} $=1{,}5$ ($30\ \mathrm{mm}$) y
\texttt{radio\_oclusion\_celdas} $=2{,}0$ ($40\ \mathrm{mm}$), medidos entre centros. La
distancia real entre centros en el instante de contacto es
\[
  \varrho_{\mathrm{cont}} = \tfrac{L}{2}\big(|\cos\Delta|+|\sin\Delta|\big) + \tfrac{\Lambda}{2}
  \in [\,100{,}0;\ 112{,}4\,]\ \mathrm{mm},
\]
es decir entre $3{,}3$ y $3{,}7$ veces el radio declarado. Con cuerpos de
$120\times140\ \mathrm{mm}$ y cubos de $60\ \mathrm{mm}$, que el centro del rover llegue
a $30\ \mathrm{mm}$ del centro del cubo exige que el rover haya atravesado el cubo.
\campo{Impacto} Doble. \textbf{(a)} Un planificador calibrado contra el simulador
aprende a acercarse demasiado y en la cancha embiste el cubo antes de llegar a su punto
de puesta. \textbf{(b)} El simulador \emph{nunca ocluye un cubo durante un empuje
legítimo} ($5{,}0 > 2{,}0$ celdas), que es justo el caso que el sistema real reporta como
el más frecuente y que degrada la detección de $1{,}05$ a $4{,}88\ \mathrm{mm}$. La
lógica de \texttt{age\_ms} queda sin ejercitar donde más importa.
\campo{Acción} Elevar a la organización. Mientras tanto, no calibrar distancias de
aproximación contra el simulador oficial; usar el simulador de N4.
\end{hallazgo}

\begin{hallazgo}{H-02}{El rango de la diferencia angular está mal documentado}
\campo{Dónde} El modelo matemático y el docstring de
\texttt{vision/detectors/rovers.py::diferencia\_angular}, que declaran
$(-180\grados,\,180\grados]$.
\campo{Qué} La forma cerrada
$\delta(\alpha,\beta)=\big((\alpha-\beta+180\grados)\bmod 360\grados\big)-180\grados$
tiene rango $[-180\grados,\,180\grados)$: $\delta=+180\grados$ exigiría que el módulo
devolviera $360\grados$, imposible. Verificado sobre $2{,}16$ millones de evaluaciones:
$-180\grados$ se alcanza, $+180\grados$ nunca.
\campo{Impacto} Bajo, pero real. Una aserción $-180 < \delta \le 180$ falla exactamente
en el caso antipodal, que es infrecuente y por eso aparece tarde. Operativamente, un
giro de exactamente $180\grados$ se resuelve siempre hacia el mismo lado, de forma
determinista.
\campo{Acción} \textbf{Corregido en el modelo} (observación 2.5, nueva). Señalar el
docstring a la organización.
\end{hallazgo}

\begin{hallazgo}{H-03}{La bandera \texttt{confiable} no protege el ángulo}
\campo{Dónde} \texttt{vision/detectors/cubos.py}: el criterio
$J_\beta \le \tau_J = 0{,}2$ celdas.
\campo{Qué} El umbral fue calibrado contra el error de \emph{posición} y funciona muy
bien para eso: sobre $436$ observaciones acumuladas, \textbf{cero falsos negativos} con
el criterio «error de posición $>10\ \mathrm{mm}$». Aplicado al \emph{ángulo} con el
umbral de $5\grados$ que impone la restricción de interfaz RI2, el mismo criterio da
$\mathbf{20}$ \textbf{falsos negativos de} $\mathbf{35}$ casos malos: una tasa del
$\mathbf{57\%}$.

\medskip
\centering
\begin{tabular}{lrr}
\toprule
& error $>5\grados$ & error $\le5\grados$ \\
\midrule
marcado \textsc{no} confiable & $15$ & $0$ \\
marcado confiable & $\mathbf{20}$ & $401$ \\
\bottomrule
\end{tabular}

\raggedright
\medskip
La causa es la medida en N2.3: con un rover empujando ($\approx22\%$ de oclusión) la
posición se degrada $9{,}5\times$ pero el ángulo se degrada $28\times$, y el residuo
---que mide ajuste global de la silueta--- no distingue una cosa de la otra.
\campo{Impacto} Alto \emph{si algún día se publica el ángulo}. Un consumidor que
reutilice \texttt{confiable} para filtrar el ángulo aceptará errores de hasta
$19\grados$ como buenos, cuatro veces el límite de RI2.
\campo{Acción} Elevar a la organización: el ángulo necesita su propio umbral de
confiabilidad, no puede compartir el de la posición. Confirma además, por otra vía, la
decisión del modelo de planificar sin depender de $\theta$.
\end{hallazgo}

\begin{hallazgo}{H-04}{La regla de decisión avance/retroceso del protocolo es incorrecta}
\campo{Dónde} Protocolo de pruebas, N3.4, que proponía «el retroceso gana si
$|\delta(\varphi_0,\theta_0)| > 90\grados$».
\campo{Qué} Esa regla acierta solo el $\mathbf{75\%}$ de las veces sobre $10^5$ ternas
aleatorias, porque mira el primer giro e ignora el segundo. Usando la identidad
verificada $|\delta(\varphi_0+180\grados, x)| = 180\grados - |\delta(\varphi_0,x)|$
(error máximo $1{,}1\times10^{-13}$), la diferencia de costos es
\[
  T^{+} - T^{-} \;=\; \frac{2\,(A + B - 180\grados)}{\omega},
  \qquad A = |\delta(\varphi_0,\theta_0)|,\quad B = |\delta(\varphi^\star,\varphi_0)|,
\]
de donde la regla exacta es $A+B>180\grados$. Acierta el $\mathbf{100{,}00\%}$.
\campo{Impacto} Medio. La regla equivocada produce un cuarto de las maniobras con un
giro de más. Se manifiesta como «el rover da vueltas al pedo», síntoma que suele
atribuirse erróneamente al controlador de bajo nivel.
\campo{Acción} \textbf{Implementada} en el planificador de N4. Conviene agregarla al
modelo como corolario del teorema \textsc{gtg}.
\end{hallazgo}

\begin{hallazgo}{H-05}{El empuje multi-tramo es obligatorio, no opcional}
\campo{Dónde} Modelo matemático, corolario de la zona de puesta inviable, que enunciaba
la región de forma cualitativa.
\campo{Qué} Cuantificada sobre una malla de $173\times173$ posiciones de cubo para cada
depósito: el $\mathbf{36{,}3\%}$ de la cancha no admite empuje directo (modelo de
rectángulo orientado). Con el disco envolvente, más conservador, sube a $39{,}96\%$: la
diferencia de $3{,}7$ puntos es el precio de aproximar el rover por un disco.
\campo{Impacto} Medio-alto. Más de un tercio de las posiciones posibles de un cubo exige
la descomposición multi-tramo. Tratarla como una rama excepcional es un error de diseño;
hay que implementarla y probarla como camino principal.
\campo{Acción} Cuantificar en el modelo. Implementada en el planificador de N4.
\end{hallazgo}

\begin{hallazgo}{H-06}{Las poses iniciales oficiales hacen que los rovers se solapen}
\campo{Dónde} \texttt{contrato/config\_simulador.json}: rovers en $(4;\,4)$ y $(4;\,8)$.
\campo{Qué} Están a $4$ celdas $=80\ \mathrm{mm}$ entre centros, pero cada rover mide
$120\ \mathrm{mm}$ de ancho. Se solapan $40\ \mathrm{mm}$. Las separaciones mínimas son:
\[
  \text{lado a lado } 120\ \mathrm{mm}\ (6{,}0\ \text{c}),\qquad
  \text{en fila } 140\ \mathrm{mm}\ (7{,}0\ \text{c}),\qquad
  \text{cualquier rumbo } 184{,}4\ \mathrm{mm}\ (9{,}22\ \text{c}).
\]
El simulador oficial no lo detecta porque no tiene modelo de cuerpo: sus rovers son
puntos.
\campo{Impacto} Alto el día del montaje. Si la zona de salida real obliga a esa
separación, los dos robots no entran. Además, cualquier prueba que arranque desde esa
configuración empieza con los cuerpos interpenetrados.
\campo{Acción} Elevar a la organización y \textbf{medir físicamente la zona de salida}.
En el banco de N4 se usan $7$ celdas de separación.
\end{hallazgo}

\begin{hallazgo}{H-07}{La cesión por pares no basta como coordinación}
\campo{Dónde} Frontera del alcance del modelo, que excluyó explícitamente la
coordinación multi-agente.
\campo{Qué} N3.8 midió que el $\mathbf{46\%}$ de los pares de corredores de empuje
aleatorios interfiere. N4.2 midió la consecuencia en dinámica: con dos rovers y la
cesión por pares como única coordinación, la entrega se mantiene pero hay colisiones en
el $\mathbf{5{,}5\%}$ de las corridas (con las dimensiones supuestas eran el $100\%$;
el rover real, más chico, se estorba mucho menos).
\campo{Impacto} Medio, y revisado a la baja. Las colisiones no impidieron la entrega en
el simulador, pero en la cancha física una colisión puede desplazar un cubo, desalinear
un marcador o penalizar. Con $0{,}2$ colisiones de media por corrida el problema pasó de
sistemático a esporádico, pero no desapareció: la condición geométrica de N3.8 sigue
diciendo que el $46\%$ de los pares de corredores interfiere, y lo que bajó es la
probabilidad de que ese cruce se materialice, no la ausencia de coordinación.
\campo{Acción} Implementar coordinación real: reserva espacio-temporal de corredores, o
asignación que exija corredores disjuntos (la condición geométrica ya está verificada en
N3.8, con acuerdo del $99{,}99\%$ contra muestreo denso).
\end{hallazgo}

\begin{hallazgo}{H-09}{La zona trampa: un cubo pegado a la pared es irrecuperable}
\campo{Dónde} Consecuencia geométrica no prevista en el modelo, descubierta al
investigar por qué fallaban los escenarios adversarios (a) y (f) de N4.5.
\campo{Qué} Para empujar un cubo \emph{alejándolo} de una pared, el rover tiene que
colocarse entre el cubo y esa pared. Necesita $\varrho_{\mathrm{puesta}} = 7{,}62$
celdas, y el tablero físico solo ofrece $3{,}5$ celdas de borde muerto. Por lo tanto un
cubo a menos de
\[
  \varrho_{\mathrm{puesta}} - 3{,}5 = 4{,}12\ \text{celdas} = 82\ \mathrm{mm}
\]
de una pared \textbf{no puede ser alejado de ella}. Barriendo toda la cancha, el
$\mathbf{11{,}22\%}$ de las posiciones es zona trampa: ninguna dirección de empuje con
punto de puesta válido incrementa la distancia a la pared. Con la distancia de contacto
mínima (sin margen de seguridad) la zona baja solo a $9{,}53\%$: \emph{no es un problema
de margen, es un problema de tamaño del rover}.
\campo{Impacto} \textbf{Estratégico.} Un cubo empujado hacia una pared se pierde para el
resto de la ronda. Y como los depósitos están a $2{,}5$ celdas del borde, es decir
\emph{dentro} de la zona trampa, empujar un cubo hacia el depósito equivocado lo pierde
de forma definitiva.
\campo{Acción} Tres reglas para el planificador:
\begin{enumerate}[label=(\alph*),leftmargin=2em,itemsep=1pt,topsep=2pt]
\item \textbf{Nunca empujar un cubo hacia una pared} si existe alternativa; preferir
      trayectorias que lo mantengan a más de $4{,}12$ celdas del borde.
\item \textbf{Detectar la zona trampa y abandonar} el objetivo en lugar de gastar la
      ronda. Implementado: en N4.5 el planificador identifica correctamente los cubos de
      (a) y (f) como irrecuperables y libera el rover.
\item \textbf{Verificar el corredor de empuje} antes de comprometerse: si el segmento
      $[\mathbf c,\mathbf g]$ pasa cerca de una pared, hay riesgo de que una desviación
      lateral meta el cubo en la trampa.
\end{enumerate}
\end{hallazgo}

\begin{hallazgo}{H-08}{El umbral de latencia del contrato es inerte}
\campo{Dónde} Sección 6.4 del contrato oficial, adoptada por el modelo como restricción
de interfaz RI4: «frenar si la latencia supera $500\ \mathrm{ms}$».
\campo{Qué} El barrido de N4.4 muestra que el sistema deja de entregar a partir de
$200\ \mathrm{ms}$, y que la regla nunca dispara por debajo de $500$. Las columnas «con
corte» y «sin corte» son \textbf{idénticas en los siete puntos del barrido}: por debajo
del umbral la regla no actúa, y por encima actúa en un régimen donde la tasa de entrega
ya es cero. Las $14\,220$ frenadas contabilizadas a $800\ \mathrm{ms}$ no cambiaron ni
un resultado.
\campo{Impacto} Medio. La regla, tal como está especificada, no modifica el
comportamiento del sistema en ningún punto del rango probado. Da una sensación de
protección que no corresponde a ningún efecto medible.
\campo{Acción} Dos cosas, y la segunda importa más que la primera.
\begin{enumerate}[label=(\alph*),leftmargin=2em,itemsep=1pt,topsep=2pt]
\item Bajar el umbral a la zona donde todavía hay algo que proteger: la degradación es
      medible a $100\ \mathrm{ms}$ y severa a $200$.
\item \textbf{Diseñar el bucle completo para operar por debajo de $100\ \mathrm{ms}$},
      no de $500$. El presupuesto de latencia real ---captura, proceso, red, decisión,
      transmisión al rover--- es cinco veces más ajustado de lo que sugiere el contrato.
      Con la telemetría a $20\ \mathrm{Hz}$ el intervalo entre mensajes ya es
      $50\ \mathrm{ms}$: queda muy poco margen para todo lo demás.
\end{enumerate}
\end{hallazgo}

% =========================================================================
\section{Correcciones aplicadas}

\subsection{Al modelo matemático}

\begin{enumerate}[label=\arabic*.,leftmargin=2.4em,itemsep=3pt]
\item \textbf{Rango de $\delta$} (H-02): corregido a $[-180\grados,180\grados)$, con una
      observación nueva que demuestra por qué $+180\grados$ es inalcanzable y qué
      significa operativamente. El documento del modelo pasa de 27 a 28 páginas.
\item \textbf{Regla avance/retroceso} (H-04): la regla exacta $A+B>180\grados$ se deriva
      de la identidad $|\delta(\varphi_0+180\grados,x)| = 180\grados-|\delta(\varphi_0,x)|$
      y debe agregarse como corolario del teorema \textsc{gtg}.
\item \textbf{Costo de un quiebre} (N3.6): el modelo declaraba $279\ \mathrm{mm}$
      equivalentes, que es el término retirada $+$ arco. El costo completo, incluyendo la
      reorientación, es $399\ \mathrm{mm}$.
\item \textbf{Zona trampa} (H-09): resultado geométrico nuevo, no previsto en el modelo,
      que merece su propia proposición.
\end{enumerate}

\subsection{Al banco de pruebas}

\begin{enumerate}[label=\arabic*.,leftmargin=2.4em,itemsep=3pt]
\item \textbf{Modelo de oclusión}: de huellas a siluetas (ver 4.5.1). Sin esta
      corrección, el nivel N4 habría reproducido el mismo defecto que denuncia en H-01.
\item \textbf{Poses iniciales} (H-06): separadas $7$ celdas en vez de $3$.
\item \textbf{Manejo de frontera}: proyección de vuelta al tablero en lugar de reversión
      del paso, que dejaba al rover atrapado si llegaba fuera.
\item \textbf{Criterio de N3.1}: el umbral original del $35\%$ era arbitrario y sin
      fundamento. Se reemplazó por una verificación geométrica (la región debe estar del
      lado opuesto al depósito y ser coherente punto a punto), y la fracción pasó de
      criterio a resultado. Es un caso de la regla de detención (c) del protocolo: el
      criterio de aceptación estaba mal, y queda escrito que lo estaba.
\end{enumerate}

% =========================================================================
\section{Qué quedó sin cubrir}

\begin{table}[htbp]
\centering
\small
\caption{Límites declarados de esta campaña de pruebas.}
\begin{tabular}{p{4.6cm}p{9.6cm}}
\toprule
Área & Estado \\
\midrule
Mecánica del contacto & N4.7 muestra que el modelo de rotación cuasiestático usado
($r_g^2 = L^2/6$, centro instantáneo) \textbf{no es cuantitativamente confiable}:
resulta no monótono en el desplazamiento lateral $e$ y no produce deriva lateral, que
físicamente debería existir. Solo el caso $e=0$ está verificado (deriva nula a
$1{,}7\times10^{-13}$ grados, por simetría). Cuantificar la rotación inducida exige un
modelo de fricción, que está fuera del alcance declarado. \\
\addlinespace
Controlador de bajo nivel & No se sintetizó ni se probó. El planificador de N4 usa un
control proporcional mínimo, suficiente para ejercitar la geometría pero no
representativo del comportamiento del rover real con IMU y PID. \\
\addlinespace
Fusión sensorial & El sensor de color, el ultrasónico, los infrarrojos y la IMU no
participan de ninguna prueba. El lazo cerrado usa exclusivamente la telemetría de
visión. \\
\addlinespace
Coordinación multi-agente & Solo la condición geométrica de compatibilidad (N3.8) y una
regla de cesión por pares. H-07 mide lo que falta. \\
\addlinespace
Cámara real & Todas las pruebas de visión usan el generador sintético. La calibración
con cámara física, la iluminación real y el comportamiento del plástico de los cubos
quedan sin verificar. \\
\addlinespace
Comunicación & El transporte TCP/NDJSON no se ejercitó: el simulador entrega el mensaje
en memoria. La latencia se modela como retardo puro, sin pérdidas de paquetes ni
reconexiones. \\
\bottomrule
\end{tabular}
\end{table}

% =========================================================================
\section{Recomendaciones, en orden}

\begin{enumerate}[label=\textbf{\arabic*.},leftmargin=2.6em,itemsep=6pt]
\item \textbf{Implementar la estimación del ángulo del cubo} (H-11). Ya no es una
      optimización: con $W_u = 93{,}50\ \mathrm{mm}$ hay un $22{,}8\%$ de orientaciones
      en las que el empuje ciego no tiene margen de control alguno. El contrato v1 no
      publica el ángulo aunque el detector oficial lo calcule, así que la estimación
      corre del lado del cliente, con los estimadores cerrados de N1. La regla operativa
      es apuntar el empuje a menos de $21\grados$ de la normal a una cara.

\item \textbf{Rediseñar el planificador para un cuerpo largo} (H-12, H-14, H-15). Tres
      reglas, todas ya implementadas y verificadas, que no existían: acceso en dos tramos
      con anillo de maniobra de $156\ \mathrm{mm}$; retirada recta de $55\ \mathrm{mm}$
      antes de girar; y los demás cubos como obstáculos de navegación con $92\ \mathrm{mm}$
      de despeje lateral. Sin ellas el rover deshace su propio trabajo: $4$ des-entregas
      por corrida.

\item \textbf{Rehacer la medición con calibre cuando esté el rover definitivo.} El actual
      es una versión preliminar medida con regla escolar. Los tres números a confirmar:
      canal $93{,}50$, huella $99{,}50$, alcance de paletas $55$. Verificar también la zona
      de salida (H-06: ¿entran los dos rovers?).

\item \textbf{Elevar a la organización los tres hallazgos del repositorio oficial:}
      H-01 (contacto y oclusión geométricamente imposibles en el simulador), H-06 (poses
      iniciales solapadas) y H-03 (la bandera de confiabilidad no sirve para el ángulo).
      Los tres son verificables en cinco minutos y los tres afectan a todos los equipos.

\item \textbf{Diseñar el bucle para $100\ \mathrm{ms}$, no para $500$} (H-08). El
      presupuesto de latencia real es cinco veces más ajustado que el que sugiere el
      contrato.

\item \textbf{Implementar la regla de la zona trampa} (H-09) antes que cualquier
      optimización de rutas. Perder un cubo contra una pared cuesta más que cualquier
      ganancia de eficiencia, y la regla es barata: no empujar hacia el borde, y
      abandonar lo irrecuperable.

\item \textbf{Tratar el multi-tramo como camino principal} (H-05): lo necesita el
      $25{,}7\%$ de las posiciones iniciales posibles.

\item \textbf{Construir la coordinación real} (H-07), usando la condición de corredores
      disjuntos de N3.8 como criterio de asignación.

\item \textbf{No usar B4 solo} (N1.11). La cascada correcta es: B2 sobre cada arista
      candidata, mediana circular, descarte de incompatibles, y recién entonces B4 sobre
      los supervivientes.

\item \textbf{Repetir N2 con la cámara real} cuando esté montada. Es la única prueba de
      esta campaña cuyo resultado podría cambiar con hardware físico.
\end{enumerate}

\begin{tcolorbox}[colback=black!3,colframe=black!55,boxrule=0.6pt,arc=1pt]
\textbf{Conclusión.} El núcleo geométrico del modelo ---la parametrización del cuadrado,
los estimadores de pose, la cota de anchura proyectada y el embudo de tolerancia--- se
sostiene sin una sola excepción sobre millones de casos. Lo que las pruebas encontraron
no fueron errores de álgebra sino \emph{consecuencias no previstas} de esa misma
geometría: la zona trampa, la magnitud real de la región sin empuje directo y el costo
verdadero de un quiebre. Y encontraron tres defectos en la infraestructura oficial que
ningún equipo puede permitirse descubrir el día de la competencia.

\medskip
Pero la lección más cara del trabajo no la dio la matemática sino un dato: las dimensiones
del rover. Se supuso un ancho de $120\ \mathrm{mm}$ y era $98{,}66$; se creyó que $98{,}66$
era el ancho útil y era $93{,}50$; se creyó que el rover medía $94\ \mathrm{mm}$ de largo
y mide $150$. Cada corrección de esos números redujo el margen: la del ancho llevó la
tolerancia lateral por debajo de la incertidumbre de la propia visión (H-11), y la del
largo hizo crecer el envolvente un $66\%$ (H-12), con lo que el rover pasó a estorbarse a
sí mismo. La conclusión que el modelo sostenía en su primera versión ---que se podía
empujar sin conocer el ángulo del cubo--- \textbf{no sobrevivió al dato real}. Nada en el
álgebra lo anticipaba; solo medirlo lo mostró.

\medskip
Conviene enunciar el corolario metodológico, porque se repitió tres veces en este trabajo
y siempre con el mismo signo: \textbf{cada vez que un dato supuesto fue reemplazado por
uno real, el margen se achicó}. Nunca al revés. Un modelo alimentado con valores de
relleno no es un modelo optimista por casualidad: es optimista porque los valores de
relleno se eligen cómodos. La única defensa es medir temprano, y medir lo que el modelo
usa, no lo que es fácil de medir.
\end{tcolorbox}

\newpage
% =========================================================================
\section{Adenda: revisión con las dimensiones reales}
\label{sec:adenda}

Después de cerrar la campaña de pruebas surgieron cuatro correcciones de hecho que
obligan a revisar parte de los resultados. Se documentan aquí, con lo que cambió y lo que
no, en vez de reescribir el informe en silencio. La última ---el archivo de piezas
separadas--- es la que produjo el hallazgo de mayor severidad de todo el trabajo.

\subsection{Corrección 1 --- El área de juego no tiene paredes}

La superficie de $1\times1\ \mathrm{m}$ \textbf{no tiene bordes verticales}: lo que la
sobrepasa cae al vacío. El hallazgo H-09 se había formulado como «el rover choca contra
una pared y no puede ponerse detrás del cubo». El mecanismo es otro, y la consecuencia
tiene dos mitades opuestas.

\paragraph{Se relaja la restricción cinemática.} Contra una pared el rover no puede
sobresalir; contra un precipicio puede voladizo, y cae recién cuando su polígono de
apoyo deja de contener al centro de masa. La condición operativa pasa a ser que la
línea de ruedas quede sobre el tablero, no todo el cuerpo.

\paragraph{Y aparece un modo de falla terminal.} Con pared, empujar un cubo contra el
borde lo deja trabado; sin pared, el cubo \textbf{cae y se pierde}. Con pared, conducir
contra el borde produce un roce; sin pared, el rover \textbf{cae y la ronda termina}.

\subsection{Corrección 2 --- Las dimensiones estaban en el repositorio}

El ancho útil entre paletas se había señalado como el parámetro crítico a medir sobre el
robot físico. No hacía falta: el repositorio incluye
\texttt{archivos\_fabricacion/CENFOBOT\_Rover\_Rev1.dxf}, el CAD de corte de la cubierta.

\begin{tcolorbox}[colback=black!3,colframe=black!50,boxrule=0.5pt,arc=1pt]
\textbf{Escala confirmada por tres vías independientes.} El DXF del rover no declara
unidades. Se verificó contra: (a) las ranuras de ensamble, de exactamente $3{,}00$
unidades, que es el espesor del acrílico declarado en el \texttt{README}; (b) los dos
agujeros de la placa frontal, de $16{,}17$ y $16{,}12$ unidades separados $24{,}00$,
que es el patrón del ultrasónico HC-SR04 ($16\ \mathrm{mm}$ a $24\ \mathrm{mm}$); y
(c) \texttt{cubos.dxf}, que sí declara milímetros y mide $180{,}30$ de lado, los tres
cubos de $60\ \mathrm{mm}$ del reto. Las tres fijan $1$ unidad $= 1\ \mathrm{mm}$.
\end{tcolorbox}

\begin{table}[htbp]
\centering
\small
\caption{Piezas medidas sobre el DXF, y el efecto de la corrección.}
\begin{tabular}{lrrr}
\toprule
Magnitud & supuesto & real (DXF) & cambio \\
\midrule
Ancho del chasis $W$ & $120{,}00$ & $\mathbf{98{,}66}$ & $-17{,}8\%$ \\
Largo del chasis $\Lambda$ & $140{,}00$ & $\mathbf{93{,}20}$ & $-33{,}4\%$ \\
Radio envolvente $R_{\mathrm{rov}}$ & $92{,}20$ & $67{,}86$ & $-26{,}4\%$ \\
Distancia de contacto (máx.) & $112{,}43$ & $89{,}03$ & $-20{,}8\%$ \\
Distancia de puesta & $152{,}43$ & $129{,}03$ & $-15{,}4\%$ \\
Tolerancia lateral $e_{\max}$ (peor caso) & $17{,}57$ & $\mathbf{6{,}90}$ & $\mathbf{-60{,}7\%}$ \\
\bottomrule
\end{tabular}
\end{table}

El rover real es \textbf{más chico y más ágil} pero \textbf{sujeta el cubo con mucho
menos margen}.

\paragraph{Y las paletas son un canal paralelo.} Las dos paletas son paralelas entre sí,
como los rieles de una vía: una U de esquinas rectas, sin convergencia ni divergencia.
El modelo de canal del teorema del embudo describe entonces el sistema \emph{sin
aproximación}, y los valores de $e_{\max}$ no son conservadores sino exactos.

\begin{tcolorbox}[colback=black!5,colframe=black!65,boxrule=0.7pt,arc=1pt]
\textbf{Una lectura errónea de la fotografía, registrada.} En una revisión intermedia se
modelaron las paletas como un embudo en V, a partir de la vista frontal del robot armado,
donde \emph{parecen} abrirse. Es un artefacto de perspectiva: las paletas se proyectan
hacia adelante y hacia abajo, sus extremos quedan más cerca de la cámara, y la
perspectiva agranda lo cercano.

Es el segundo intento fallido de sacar geometría de esa misma fotografía: el primero
había sido escalarla, con dos estimaciones que diferían un $52\%$. La conclusión
metodológica vale más que el error: \textbf{una fotografía en perspectiva, sin
referencias métricas en el plano de interés, no sirve para deducir geometría}, ni
cuantitativa ni cualitativamente. La única fuente confiable sin el robot en la mano es
el archivo de fabricación.

\medskip
Ningún número de este informe cambia por la corrección: el código del banco siempre
implementó el canal paralelo. El embudo existió solo en la prosa de esa revisión.
\end{tcolorbox}

\paragraph{Y el cubo puede girar libremente dentro del canal.} Como
$W_u = 93{,}50 > 84{,}85 = L\sqrt2$, no hay orientación que trabe el cubo entre las
paletas. Eso descarta el atasco, pero también significa que \textbf{nada impide que un
empuje descentrado lo haga rotar hasta $\Delta = 45\grados$}, que es la orientación de
peor tolerancia. Un cubo que entra alineado y gira durante el transporte ve su holgura
caer de $16{,}75$ a $4{,}32\ \mathrm{mm}$ sin que el sistema lo advierta.

\subsection{Corrección 3 --- El archivo de piezas separadas cierra la geometría}

El ancho útil del canal $W_u$ había quedado acotado a $[92{,}66;\ 98{,}66]\ \mathrm{mm}$
---según las paletas montaran por dentro o por fuera de los paneles laterales--- y
señalado como la única medición pendiente. El archivo
\texttt{CENFOBOT\_PIEZAS\_SEPARADAS.dxf}, que contiene las piezas de fabricación
desplegadas una junto a otra, la resuelve por construcción.

\begin{tcolorbox}[colback=black!3,colframe=black!50,boxrule=0.5pt,arc=1pt]
\textbf{Cómo se resuelve, sin ambigüedad.} El chasis mide $93{,}50 \times 94{,}00$ de
cuerpo, y sus bordes verticales están en $224{,}85,\ 227{,}85,\ \dots,\ 321{,}35,\ 324{,}35$:
las lengüetas sobresalen exactamente $3{,}00$ por lado, que es el espesor del acrílico.
Una lengüeta de espesor igual al del material es una \textbf{junta pasante}: atraviesa la
ranura del panel lateral y queda enrasada con su cara exterior. Por lo tanto la cara
\emph{interior} del panel coincide con el borde del cuerpo del chasis, y
\[
  W_u = 93{,}50\ \mathrm{mm}, \qquad
  W_{\mathrm{ext}} = 93{,}50 + 2\cdot 3{,}00 = 99{,}50\ \mathrm{mm}.
\]
Esto confirma exactamente la fórmula que había propuesto el autor: ancho total menos los
$6\ \mathrm{mm}$ de los dos espesores. La medición con calibre sobre el robot armado
sigue siendo deseable como verificación, pero ya no es una incógnita del modelo.
\end{tcolorbox}

\begin{table}[htbp]
\centering
\small
\caption{Piezas del archivo de fabricación separado. Escala: $1$ unidad $= 1\ \mathrm{mm}$
(\texttt{\$INSUNITS} $=4$, milímetros, declarado en el archivo).}
\begin{tabular}{lrl}
\toprule
Pieza & Dimensiones (mm) & Nota \\
\midrule
Chasis (cuerpo) & $93{,}50 \times 94{,}00$ & define $W_u$ y $\Lambda$ \\
Chasis (con lengüetas) & $99{,}50 \times 94{,}00$ & huella exterior $W_{\mathrm{ext}}$ \\
Panel lateral (${\times}2$) & $52{,}50 \times 106{,}00$ & \\
Placa frontal & $99{,}38 \times 42{,}63$ & tramos rectos $99{,}00$ \\
Tapa superior & $72{,}09 \times 51{,}53$ & \\
Placa trasera & $66{,}00 \times 22{,}00$ & \\
Soporte de motor (${\times}2$) & $47{,}90 \times 28{,}60$ & \\
Lengüeta pasante & $3{,}00$ por lado & $=$ espesor del acrílico \\
\bottomrule
\end{tabular}
\end{table}

\begin{tcolorbox}[colback=black!5,colframe=black!65,boxrule=0.7pt,arc=1pt]
\textbf{Dos anchos distintos, y confundirlos es un error.} El modelo usa ahora
\emph{dos} anchos que no son intercambiables:
\begin{itemize}[leftmargin=1.6em,itemsep=2pt,topsep=3pt]
\item $W_{\mathrm{ext}} = 99{,}50\ \mathrm{mm}$, la \textbf{huella exterior}. Gobierna
      colisiones, radio envolvente $R_{\mathrm{rov}}$ e inflado de obstáculos.
\item $W_u = 93{,}50\ \mathrm{mm}$, el \textbf{ancho útil del canal}. Gobierna
      $e_{\max}$ y todo el presupuesto de error lateral.
\end{itemize}
Usar $W_{\mathrm{ext}}$ donde va $W_u$ infla la tolerancia lateral un $69\%$
($4{,}32 \to 7{,}32\ \mathrm{mm}$ en el peor caso) y hace desaparecer el hallazgo H-11.
El banco de pruebas fue corregido para llevar las dos constantes por separado.
\end{tcolorbox}

\subsection{Qué cambió en los resultados}

\begin{table}[htbp]
\centering
\small
\caption{Resultados revisados. Todo lo que no aparece aquí no cambió.}
\begin{tabular}{llrr}
\toprule
Prueba & Magnitud & antes & después \\
\midrule
N3.1 & Cancha sin empuje directo & $36{,}3\%$ & $25{,}7\%$ \\
N3.2 & $e_{\max}$ & $[17{,}57;\ 30{,}00]$ & $[\mathbf{4{,}32};\ 16{,}75]$ mm \\
N3.5 & Umbral de cuerda del anillo & $84{,}9\grados$ & $92{,}59\grados$ \\
N3.6 & Costo de un quiebre de $90\grados$ & $399\ \mathrm{mm}$ & $363\ \mathrm{mm}$ \\
N3.9 & Zona de exclusión (pared) & $11{,}22\%$ & $5{,}03\%$ \\
N3.9 & Zona de exclusión (precipicio) & --- & $\mathbf{2{,}76\%}$ \\
N3.10 & Ventana angular admisible & --- & $|\Delta| \le 20{,}97\grados$ \\
N3.10 & Banda de presupuesto negativo & --- & $\mathbf{(34{,}76;\ 55{,}24)\grados}$ \\
N4.1 & Misión nominal, 1 rover & $3/3$ en $33{,}0$ s & $3/3$ en $33{,}4$ s \\
N4.1 & Misión nominal, 2 rovers & $3/3$, 2 colisiones & $3/3$, \textbf{0 colisiones} \\
N4.2 & Tasa de entrega, 1 rover & $99{,}5\%$ & $98{,}5\%$ \\
N4.2 & Corridas limpias, 2 rovers & $0\%$ & $\mathbf{94{,}5\%}$ \\
\bottomrule
\end{tabular}
\end{table}

Las diez pruebas de N3 ---la nueva N3.10 incluida--- vuelven a pasar sin excepción con
los parámetros nuevos, igual que N0 a N2.
El rover más chico \emph{mejora} el desempeño en lazo cerrado: llega a más lugares,
necesita menos espacio detrás del cubo y estorba menos al otro rover.

\subsection{Lo que empeoró, y es lo que importa}

\begin{tcolorbox}[colback=black!5,colframe=black!65,boxrule=0.7pt,arc=1pt]
\textbf{El presupuesto de control dejó de existir en el peor caso.} Con
$e_{\max} = 4{,}32\ \mathrm{mm}$ y $\sigma_{\mathrm{vis}} = 5$, el presupuesto de control
es
\[
  4{,}32 - 5{,}00 \;=\; \mathbf{-0{,}68\ \mathrm{mm}} ,
\]
negativo. No hay precisión de rumbo que rescate ese caso: la incertidumbre de la
medición, sola, ya excede la tolerancia. La banda afectada es
$\Delta \in (34{,}76\grados;\ 55{,}24\grados)$, el $22{,}8\%$ de las orientaciones.

Con el empuje alineado a una cara ($\Delta \equiv 0 \bmod 90\grados$), $e_{\max}$ sube a
$16{,}75\ \mathrm{mm}$, el presupuesto a $11{,}75$, y la exigencia de rumbo pasa a ser
\[
  \epsilon_\theta \le \arcsin\!\left(\frac{11{,}75}{129{,}43}\right) = \mathbf{5{,}21\grados},
\]
que sí es alcanzable. La ventana práctica, reservando $3\ \mathrm{mm}$ de margen de
control, es $|\Delta| \le 20{,}97\grados$: el $46{,}6\%$ de las orientaciones, y cuatro
direcciones admisibles por cubo.

\medskip
\textbf{Esto no matiza la conclusión número dos del modelo original: la invierte.} Con
el ancho supuesto, conocer $\theta$ era una mejora de eficiencia. Con el cuerpo del
chasis, una palanca de robustez. Con el ancho útil real del canal es \textbf{condición
de funcionamiento}: hay orientaciones en las que el empuje ciego no tiene margen alguno.
El error angular del detector oficial con el cubo despejado es de $0{,}43\grados$ de
media (N2.2), holgadamente suficiente para satisfacer la ventana --- pero el contrato v1
no publica el ángulo, de modo que la estimación hay que hacerla del lado del cliente.
\end{tcolorbox}

\subsection{Lo que falta medir ahora}

El archivo de piezas separadas cerró la incógnita que gobernaba el modelo. Queda muy
poco.

\begin{enumerate}[label=\textbf{\arabic*.},leftmargin=2.6em,itemsep=3pt]
\item \textbf{Verificación, no medición:} confirmar con calibre que el robot armado mide
      $99{,}50\ \mathrm{mm}$ punta a punta. Si el valor coincide, $W_u = 93{,}50$ queda
      confirmado y con él todo el hallazgo H-11.
\item $\lambda$, el alcance de las paletas por delante de la placa frontal. \textbf{No
      afecta a $e_{\max}$}, que es la magnitud crítica; fija la distancia de contacto y
      con ella el punto de puesta, y el modelo lo trata mediante una cota válida
      cualquiera sea $\lambda$.
\end{enumerate}

El ángulo de divergencia salió del problema al confirmarse que las paletas son
paralelas; el ancho útil salió al aparecer el archivo de piezas separadas. La
parametrización geométrica del modelo está cerrada.

\subsection{Corrección 4 --- El rover físico, medido}
\label{sec:medido}

Se midió el rover armado con regla escolar sin calibrar, sobre una \textbf{versión
preliminar} del robot. Las lecturas se tratan como $\pm 1\ \mathrm{mm}$ y ninguna
conclusión se apoya en un solo número.

\paragraph{Cómo se mide con regla lo que parece necesitar calibrador.} El valor en disputa
es $W_u$ y el umbral que decide es $94{,}85\ \mathrm{mm}$. Separar $93{,}5$ de $94{,}9$
sobre un tramo de $94\ \mathrm{mm}$ excede la resolución del instrumento; la misma
diferencia leída sobre un \emph{hueco de $9\ \mathrm{mm}$}, no. Todas las mediciones
críticas se plantearon como huecos.

\begin{table}[htbp]
\centering
\small
\caption{Mediciones sobre el rover armado.}
\begin{tabular}{clrrl}
\toprule
\# & Magnitud & predicho & medido & lectura \\
\midrule
M1  & Hueco con cubo a $45\grados$ & $8{,}65$ & $9{-}10$ & $W_u = 94{,}35$ \\
M2  & Hueco con cubo alineado & $33{,}50$ & $33$ & $W_u = 93{,}00$ \\
M9  & Canal, medido directo & $93{,}50$ & $95$ & $W_u = 95{,}00$ \\
\midrule
M6  & Alcance de paletas $\lambda$ & --- & $\mathbf{55}$ & nuevo \\
M7  & Ancho máximo real & $99{,}50$ & $100$ & concuerda \\
M8  & Largo máximo real & $149$ & $150$ & concuerda \\
M10 & Paletas: piso a borde superior & --- & $0$ a $14$ & nuevo \\
M12 & Aristas del cubo & vivas & vivas & concuerda \\
M13 & Rueda: diámetro / trocha & --- & $33$ / $89$ & nuevo \\
\bottomrule
\end{tabular}
\end{table}

\begin{tcolorbox}[colback=black!3,colframe=black!50,boxrule=0.5pt,arc=1pt]
\textbf{El control interno se cumple.} La resta $\text{M2}-\text{M1} = 23{,}5\ \mathrm{mm}$
contra los $L(\sqrt2-1) = 24{,}85$ que exige la identidad geométrica: desvío de
$1{,}35\ \mathrm{mm}$, compatible con dos lecturas de regla. Las mediciones son
internamente coherentes.
\end{tcolorbox}

\paragraph{H-11 queda confirmado, y de forma robusta al error.} Las tres estimaciones de
$W_u$ ---$94{,}35$, $93{,}00$ y $95{,}00$--- promedian $94{,}05 \pm 0{,}54\ \mathrm{mm}$.
El valor deducido del archivo de fabricación, $93{,}50$, cae a $1{,}0\sigma$: la medición
\emph{no lo refuta}. Y lo que importa no es el valor sino su consecuencia:
\[
  e^{\max}_{\mathrm{control}}(45\grados) \in [\,-0{,}94;\ +0{,}13\,]\ \mathrm{mm}
  \quad \text{para todo } W_u \text{ en } \pm2\sigma .
\]
\textbf{El presupuesto lateral es cero dentro del error de medición, cualquiera sea el
valor que se tome.} No hace falta un instrumento mejor para sostener la conclusión.

\paragraph{Lo que la medición sí refutó.} El alcance de las paletas,
$\lambda = 55\ \mathrm{mm}$, no estaba en el archivo de fabricación y el modelo lo había
esquivado con una cota. Medido, produce el hallazgo H-12: el envolvente real es
$113{,}49\ \mathrm{mm}$ y no $68{,}44$. De ahí se desprendieron, al rehacer el simulador,
los hallazgos H-13, H-14 y H-15.

\paragraph{Lo que confirmó, y el modelo no sabía.} Las paletas llegan al piso sin
separación medible: el polígono de apoyo son las dos ruedas más las puntas de los brazos,
exactamente como se había postulado \emph{sin evidencia} en el modelo del precipicio. De
ahí se sigue que una paleta en voladizo sobre el vacío no vuelca al rover, y por eso la
restricción de frontera se evalúa sobre el chasis y no sobre la huella completa. Además,
las paletas abarcan de $0$ a $14\ \mathrm{mm}$ de altura, así que el empuje actúa muy por
debajo de la mitad del cubo: no hay riesgo de vuelco y la hipótesis H6 queda respaldada.

\subsection{El simulador ahora tiene forma de U}

Hasta esta revisión el rover era \textbf{un solo rectángulo} en el simulador, de modo que
el cubo era empujado por la cara frontal del chasis y \textbf{el canal entre paletas no
existía}. Con eso, el resultado central del modelo ---el embudo de tolerancia, y con él
H-11--- no podía ser ni confirmado ni refutado en lazo cerrado: la prueba no tocaba el
mecanismo.

El rover se modela ahora como tres piezas: el cuerpo, que empuja, y dos paletas, que
contienen sin empujar ---resuelven la penetración lateral sin transmitir avance ni
rotación, que es exactamente lo que dice la proposición del canal---. Es el cambio que
destapó H-13, H-14 y H-15.

\subsection{El resultado que separa la geometría del planificador}

\begin{tcolorbox}[colback=black!5,colframe=black!65,boxrule=0.7pt,arc=1pt]
\textbf{Prueba N4.8, y es la más informativa de toda la campaña.} Treinta semillas, un
rover, los mismos tres cubos:

\medskip
\centering
\begin{tabular}{lrr}
\toprule
Configuración & entregados de 3 & des-entregas \\
\midrule
Un cubo por vez, misiones independientes & $\mathbf{3{,}00}$ & $0{,}00$ \\
Los tres a la vez & $\mathbf{1{,}00}$ & $\mathbf{4{,}00}$ \\
\bottomrule
\end{tabular}

\raggedright
\medskip
\textbf{La geometría del empuje no es el problema.} Un cubo solo se entrega siempre, con
la geometría corregida y el canal real. Lo que colapsa ---un $67\%$ de caída--- es hacer
tres seguidos sin perturbar a los otros dos.

\smallskip
Y el mecanismo está medido: $4$ des-entregas por corrida. El rover no falla en entregar,
\textbf{deshace entregas ya conseguidas}. Es una falla del planificador, no del modelo
geométrico: la capa de navegación con obstáculos que hace falta ---inflado de Minkowski,
grafo de visibilidad, corredores disjuntos--- está desarrollada en el modelo, pero el
planificador nunca la usó porque con la huella supuesta no hacía falta.
\end{tcolorbox}

\begin{tcolorbox}[colback=black!3,colframe=black!50,boxrule=0.5pt,arc=1pt]
\textbf{Honestidad sobre qué se está midiendo.} La caída de la misión nominal de $3/3$ a
$1/3$ es un enunciado \emph{conjunto} sobre la geometría real \emph{y} este planificador
en particular, que fue escrito para un rover compacto. Un planificador rediseñado para un
cuerpo largo con paletas debería recuperar buena parte de esa pérdida: N4.8 muestra que la
geometría lo permite. Lo que no se recupera con software son H-11 y H-12, que son
consecuencias de las dimensiones y valen para cualquier planificador.
\end{tcolorbox}

\subsection{Balance de las tres correcciones}

\begin{table}[htbp]
\centering
\small
\caption{Qué aportó cada corrección de hecho.}
\begin{tabular}{p{3.4cm}p{5.4cm}p{5.4cm}}
\toprule
Corrección & Qué cambió & Qué no cambió \\
\midrule
\textbf{Sin paredes} &
La zona de exclusión baja de $11{,}22\%$ a $2{,}76\%$: el rover puede quedar en
voladizo. Aparece un modo de falla terminal: cubo o rover que caen se pierden. &
La restricción geométrica sigue existiendo. Para alejar un cubo del borde el rover tiene
que ponerse entre el cubo y el vacío. \\
\addlinespace
\textbf{Dimensiones del DXF} &
$e_{\max}$ cae un $60{,}7\%$. La precisión de rumbo exigida pasa de $4{,}73\grados$ a
$0{,}84\grados$ en el peor caso. Conocer $\theta$ deja de ser opcional. &
Los estimadores, la cota de anchura y toda la matemática de N0 a N2. El rover más chico
\emph{mejora} el desempeño en lazo cerrado. \\
\addlinespace
\textbf{Paletas paralelas} &
La lectura del modelo: $e_{\max}$ pasa de cota conservadora a valor exacto. Se agrega la
libertad de giro del cubo dentro del canal. &
\textbf{Ningún número.} El código siempre implementó el canal paralelo; el embudo estuvo
solo en la prosa. \\
\addlinespace
\textbf{DXF de piezas separadas} &
$W_u = 93{,}50$ queda determinado por la junta pasante. $e_{\max}$ cae a
$[4{,}32;\ 16{,}75]$ y el presupuesto de control se vuelve \textbf{negativo} en el
$22{,}8\%$ de las orientaciones (H-11). Aparece la ventana angular $|\Delta|\le 21\grados$
como restricción de interfaz. &
La factibilidad: $L\sqrt2 = 84{,}85 < 93{,}50$, el cubo sigue entrando siempre. El
desempeño en lazo cerrado, que no modela el error lateral fino. \\
\addlinespace
\textbf{Medición del rover físico} &
Confirma $W_u$ y con él H-11, de forma robusta al error de la regla. Mide
$\lambda = 55\ \mathrm{mm}$, que nadie tenía: el envolvente crece un $66\%$ (H-12), y
al rehacer el simulador con forma de U aparecen H-13, H-14 y H-15. La misión nominal cae
de $3/3$ a $1/3$. &
Todo N0 a N2, y las cotas geométricas: el error de la regla ($\pm1\ \mathrm{mm}$) no
alcanza para mover ninguna conclusión. Un cubo aislado se sigue entregando siempre. \\
\bottomrule
\end{tabular}
\end{table}

\begin{tcolorbox}[colback=black!3,colframe=black!55,boxrule=0.6pt,arc=1pt]
\textbf{Lo que las tres correcciones tienen en común.} Ninguna vino de una prueba: las
tres vinieron de conocimiento del hardware que el modelo no tenía. Las pruebas
verificaron que la matemática era correcta ---y lo era, sin excepción en $46$ pruebas---
pero no podían detectar que estaba describiendo un rover de $120\ \mathrm{mm}$ con
paredes y embudo, que no existe.

Es el límite estructural de toda la campaña, y conviene tenerlo escrito:
\textbf{la verificación demuestra consistencia, no correspondencia}. Para lo segundo hace
falta o el hardware, o un archivo de fabricación, o alguien que conozca el sistema. En
este caso llegaron los tres, en ese orden inverso.
\end{tcolorbox}

\end{document}
